\documentclass[11pt,twocolumn]{article}

% ============================================================
% Packages
% ============================================================
\usepackage[utf8]{inputenc}
\usepackage[T1]{fontenc}
\usepackage{lmodern}
\usepackage[margin=1in]{geometry}
\usepackage{amsmath,amssymb,amsthm}
\usepackage{algorithm}
\usepackage{algpseudocode}
\usepackage{graphicx}
\usepackage{booktabs}
\usepackage{hyperref}
\usepackage{xcolor}
\usepackage{listings}
\usepackage{multicol}
\usepackage{enumitem}
\usepackage{caption}
\usepackage{subcaption}
\usepackage{tikz}
\usetikzlibrary{arrows.meta,positioning,shapes.geometric,fit,calc}

% ============================================================
% Listings config
% ============================================================
\lstset{
  basicstyle=\ttfamily\scriptsize,
  breaklines=true,
  frame=single,
  numbers=left,
  numberstyle=\tiny\color{gray},
  keywordstyle=\color{blue!70!black}\bfseries,
  commentstyle=\color{green!50!black}\itshape,
  stringstyle=\color{red!60!black},
  showstringspaces=false,
  tabsize=2,
}

% ============================================================
% Theorem environments
% ============================================================
\newtheorem{theorem}{Theorem}[section]
\newtheorem{lemma}[theorem]{Lemma}
\newtheorem{definition}[theorem]{Definition}
\newtheorem{property}[theorem]{Property}

% ============================================================
% Title
% ============================================================
\title{
  \textbf{Lux Teleport: A Sovereign Omnichain Liquidity Protocol} \\[6pt]
  \large Non-Custodial Cross-Chain Asset Transfer via \\
  Threshold Cryptography and Per-Chain Governance
}

\author{
  Lux Industries, Inc. \\
  \texttt{research@lux.network} \\[4pt]
  \small Version 1.0 --- April 2026
}

\date{}

\begin{document}
\maketitle

% ============================================================
% Abstract
% ============================================================
\begin{abstract}
We present Lux Teleport, a non-custodial omnichain bridge and liquidity protocol that natively connects 20 execution environments and 270+ blockchains through a unified MPC threshold signature network.
Teleport achieves trustless cross-chain asset transfer through FROST and CGGMP21 threshold cryptography with on-chain signature verification, per-chain nonce tracking, and automatic undercollateralization detection.
The protocol introduces three innovations:
(1)~\emph{sovereign chain governance} where each connected chain's native token holders control their own bridge parameters;
(2)~a \emph{yield-bearing bridge token} system where locked assets earn yield on source chains, reflected in bridge token share prices on destination chains; and
(3)~a \emph{Shariah-compliant mode} that programmatically filters yield sources to enable Islamic finance-compatible DeFi.
We prove conservation, uniqueness, liveness, and safety properties of the protocol, and describe native bridge contracts across every major smart contract runtime.
\end{abstract}

% ============================================================
\section{Introduction}
% ============================================================

Cross-chain interoperability remains the most critical unsolved infrastructure problem in distributed ledger systems. The total value locked across bridges exceeds \$15B, yet bridge exploits account for over \$2.8B in losses since 2020~\cite{rekt}. Existing approaches fall into four categories:

\begin{enumerate}[nosep]
  \item \textbf{Trusted multisig}: Custodial risk, single points of failure (WBTC, most centralized bridges).
  \item \textbf{Optimistic}: Long withdrawal delays, fraud proof windows (Across, Hop).
  \item \textbf{Light client}: Chain-specific, high implementation complexity (IBC, Wormhole).
  \item \textbf{MPC threshold}: Distributed key management, no single custodian (tBTC, Teleport).
\end{enumerate}

Teleport addresses the key management complexity of MPC bridges through two complementary threshold signature protocols:

\begin{itemize}[nosep]
  \item \textbf{FROST}~\cite{frost}: Flexible Round-Optimized Schnorr Threshold signatures for Ed25519 and BIP-340 Taproot, providing 64-byte signatures with 2-round signing.
  \item \textbf{CGGMP21}~\cite{cggmp21}: Threshold ECDSA with identifiable aborts for secp256k1 chains, producing standard 65-byte ECDSA signatures.
\end{itemize}

The combination provides native signature verification on \emph{every} major blockchain without requiring chain-specific verification contracts or precompiles.

% ============================================================
\section{System Model}
% ============================================================

\subsection{Entities}

\begin{definition}[MPC Signer Set]
A set $\mathcal{S} = \{s_1, s_2, \ldots, s_n\}$ of $n$ nodes, each holding a share of a distributed secret key. A threshold $t \leq n$ defines the minimum number of nodes required to produce a valid signature. In the default configuration, $n = 3, t = 2$.
\end{definition}

\begin{definition}[OmnichainRouter]
An immutable smart contract deployed on each connected chain that:
\begin{enumerate}[nosep]
  \item Verifies MPC threshold signatures for mint operations
  \item Tracks per-source-chain nonce bitmaps for replay prevention
  \item Distributes bridge fees to sovereign stakeholders
  \item Auto-pauses on undercollateralization
\end{enumerate}
\end{definition}

\begin{definition}[Sovereign Governor]
The native chain's DAO/timelock address that controls economic parameters (fees, limits, strategy routing) but \emph{cannot} override security invariants (nonce checks, signature verification, collateralization thresholds).
\end{definition}

\subsection{Threat Model}

We assume a computationally bounded adversary $\mathcal{A}$ who may:
\begin{itemize}[nosep]
  \item Compromise up to $t - 1$ MPC signer nodes
  \item Control the governor address on any single chain
  \item Observe and reorder transactions on any chain
  \item Deny service to individual relay nodes
\end{itemize}

We assume $\mathcal{A}$ \emph{cannot}:
\begin{itemize}[nosep]
  \item Compromise $\geq t$ MPC signer nodes simultaneously
  \item Break ECDSA, Ed25519, or Schnorr assumptions
  \item Rewrite confirmed blocks on any connected chain
\end{itemize}

% ============================================================
\section{Protocol Description}
% ============================================================

\subsection{Bridge Operations}

The protocol supports four atomic operations:

\paragraph{Lock.} User deposits asset $A$ of amount $a$ on source chain $C_s$ into the YieldBridgeVault. The vault emits event $\texttt{Lock}(C_s, C_d, \nu, A, a, r)$ where $C_d$ is the destination chain, $\nu$ is a monotonically increasing nonce, and $r$ is the recipient address on $C_d$.

\paragraph{Mint.} MPC signer set $\mathcal{S}$ observes the Lock event, reaches quorum, and produces threshold signature $\sigma$ over message:
\begin{equation}
  m = H(\texttt{"DEPOSIT"} \| C_s \| \nu \| T \| r \| a)
\end{equation}
where $H$ is keccak256 (for EVM chains) or SHA-256 (for non-EVM), and $T$ is the bridge token address on $C_d$. Any relayer submits $(C_s, \nu, T, r, a, \sigma)$ to the OmnichainRouter on $C_d$, which verifies $\sigma$ and mints $a - f$ tokens to $r$, where $f$ is the bridge fee.

\paragraph{Burn.} User submits burn transaction on chain $C_d$, destroying bridge tokens and emitting $\texttt{Burn}(C_d, C_s, \nu', T, a, r')$ where $r'$ is the recipient on $C_s$.

\paragraph{Release.} MPC signer set observes the Burn event, produces threshold signature $\sigma'$ over:
\begin{equation}
  m' = H(\texttt{"RELEASE"} \| C_d \| \nu' \| T \| r' \| a)
\end{equation}
The signature is submitted to $C_s$, which releases locked assets from the vault to $r'$.

\subsection{Nonce Management}

Each OmnichainRouter maintains a per-source-chain nonce bitmap:
\begin{equation}
  \mathcal{N}: (\text{chainId}, \text{nonce}) \rightarrow \{0, 1\}
\end{equation}

Before processing any mint or release, the router checks $\mathcal{N}(C_s, \nu) = 0$ and sets $\mathcal{N}(C_s, \nu) \leftarrow 1$ atomically within the same transaction. This prevents replay attacks across all chains.

\subsection{Daily Mint Limits}

Per-token daily limits provide defense-in-depth:
\begin{equation}
  \texttt{dailyMinted}[T] + a \leq \texttt{dailyMintLimit}[T]
\end{equation}
The counter resets every 86,400 seconds (24 hours). Limits are configurable by the sovereign governor.

% ============================================================
\section{Threshold Cryptography}
% ============================================================

\subsection{FROST Protocol}

For chains supporting Ed25519 or Schnorr signatures (Solana, TON, Sui, Aptos, Cosmos, Bitcoin Taproot, Polkadot, NEAR, Stellar, Cardano, Algorand, Tezos), we employ FROST~\cite{frost}:

\begin{algorithm}
\caption{FROST Threshold Signing ($t$-of-$n$)}
\begin{algorithmic}[1]
\Require Message $m$, participant set $P \subseteq \mathcal{S}$ with $|P| \geq t$
\Ensure Threshold signature $\sigma = (R, s)$
\State \textbf{Round 1:} Each $p_i \in P$ generates nonce pair $(d_i, e_i) \gets \mathbb{Z}_q^2$
\State Compute commitment $(D_i, E_i) = (g^{d_i}, g^{e_i})$ and broadcast
\State \textbf{Round 2:} Compute binding factor $\rho_i = H_1(i, m, \{(D_j, E_j)\})$
\State Compute group commitment $R = \prod_{i \in P} D_i \cdot E_i^{\rho_i}$
\State Compute challenge $c = H_2(R, Y, m)$ where $Y$ is the group public key
\State Each $p_i$ computes partial signature $z_i = d_i + e_i \cdot \rho_i + \lambda_i \cdot s_i \cdot c$
\State Aggregate: $s = \sum_{i \in P} z_i$
\State \Return $\sigma = (R, s)$
\end{algorithmic}
\end{algorithm}

\noindent
where $\lambda_i$ are Lagrange interpolation coefficients, $s_i$ are secret key shares, and $H_1, H_2$ are domain-separated hash functions.

\begin{theorem}[FROST Security~\cite{frost}]
Under the Discrete Logarithm assumption in the Random Oracle Model, FROST is existentially unforgeable under chosen-message attack for any adversary corrupting at most $t - 1$ of $n$ participants.
\end{theorem}

\subsection{CGGMP21 Protocol}

For EVM chains and TRON (secp256k1 ECDSA), we employ CGGMP21~\cite{cggmp21}:

\begin{enumerate}[nosep]
  \item \textbf{Key Generation}: Distributed key generation produces shares $x_i$ such that $\sum \lambda_i x_i = x$ (secret key) and $X = g^x$ (public key).
  \item \textbf{Pre-Signing}: Generates nonce shares and auxiliary ZK proofs (can be done offline).
  \item \textbf{Signing}: Given message hash $m$, produces standard ECDSA signature $(r, s, v)$ where $r = R_x \mod q$ and $s = k^{-1}(m + r \cdot x) \mod q$.
  \item \textbf{Identifiable Abort}: If a participant deviates from the protocol, the identity of the malicious party is revealed.
\end{enumerate}

\begin{theorem}[CGGMP21 Security~\cite{cggmp21}]
The CGGMP21 protocol realizes the ideal ECDSA threshold signing functionality in the presence of a malicious adversary corrupting at most $t - 1$ of $n$ parties, with identifiable abort.
\end{theorem}

\subsection{Taproot Integration}

For Bitcoin L1 custody, we use FROST in Taproot mode (BIP-340/341~\cite{bip340,bip341}):

\begin{enumerate}[nosep]
  \item \textbf{Key Generation}: $\texttt{frost.KeygenTaproot}()$ produces an $x$-only public key $\tilde{X}$ (32 bytes).
  \item \textbf{Address Derivation}: The Taproot address is $\texttt{bc1p}\langle\text{Bech32m}(\tilde{X})\rangle$.
  \item \textbf{Spending}: A Taproot spend requires a valid BIP-340 Schnorr signature, produced by the 2-of-3 FROST signing protocol.
\end{enumerate}

No on-chain smart contract is required. The MPC-controlled Taproot address \emph{is} the bridge vault. Bitcoin Script's native Schnorr verification validates the threshold signature.

% ============================================================
\section{Sovereign Governance}
% ============================================================

\subsection{Separation of Concerns}

Teleport cleanly separates three authority domains:

\begin{table}[h]
\centering
\small
\begin{tabular}{lll}
\toprule
\textbf{Domain} & \textbf{Controller} & \textbf{Scope} \\
\midrule
Security & MPC signer set & Mint/release authorization \\
Economics & Chain governor & Fees, limits, strategies \\
Protocol & Immutable code & Nonce checks, verification \\
\bottomrule
\end{tabular}
\caption{Authority domain separation}
\label{tab:separation}
\end{table}

\noindent
Critically, no single domain can override another:
\begin{itemize}[nosep]
  \item The governor \emph{cannot} mint tokens (requires MPC signature).
  \item The MPC set \emph{cannot} change fees (requires governor).
  \item Neither can bypass nonce checks (immutable code).
\end{itemize}

\subsection{Governor Capabilities}

The sovereign governor (typically a DAO timelock) can:
\begin{itemize}[nosep]
  \item Set bridge fee rate: $0 \leq f_b \leq 100$ bps (hard-capped at 1\%)
  \item Set stakeholder share: $0 \leq f_s \leq 10{,}000$ bps
  \item Register new bridge tokens
  \item Set daily mint limits per token
  \item Enable Shariah compliance mode
  \item Change fee recipient (48-hour timelock)
  \item Transfer governance (2-step process)
\end{itemize}

\subsection{Signer Rotation}

MPC signer rotation is timelocked to protect users:

\begin{enumerate}[nosep]
  \item Current 2-of-3 MPC set signs rotation proposal.
  \item 7-day timelock begins. Pending signers are visible on-chain.
  \item During timelock, users can exit (burn bridge tokens, receive underlying).
  \item After 7 days, anyone can execute the rotation.
  \item Current signers can cancel during the timelock.
\end{enumerate}

% ============================================================
\section{Yield-Bearing Bridge Tokens}
% ============================================================

\subsection{Share-Based Accounting}

Yield-bearing bridge tokens (yLETH, yLBTC, yLUSD) use ERC-4626-like share accounting:

\begin{equation}
  \text{sharePrice} = \frac{\text{totalBacking}_{\text{source}}}{\text{totalShares}_{\text{dest}}}
\end{equation}

When yield accrues on the source chain (via Lido, EigenLayer, Babylon, etc.), the MPC periodically attests to the new total backing. The share price increases automatically.

\subsection{Yield Strategy Interface}

Each strategy implements a minimal interface:

\begin{lstlisting}[language=Solidity,caption=IYieldStrategy interface]
interface IYieldStrategy {
    function deposit(uint256 amount)
        external returns (uint256 shares);
    function withdraw(uint256 shares)
        external returns (uint256 assets);
    function totalAssets()
        external view returns (uint256);
    function currentAPY()
        external view returns (uint256);
    function harvest()
        external returns (uint256);
    function isActive()
        external view returns (bool);
}
\end{lstlisting}

\subsection{Strategy Taxonomy}

\begin{table}[h]
\centering
\small
\begin{tabular}{llr}
\toprule
\textbf{Category} & \textbf{Protocols} & \textbf{APY} \\
\midrule
Liquid Staking & Lido, Rocket Pool & 3--5\% \\
Restaking & EigenLayer, Symbiotic, Karak & 3--8\% \\
Lending & Aave, Compound, Morpho & 2--12\% \\
Stablecoin & MakerDAO/Sky, Ethena & 4--15\% \\
LP/DEX & Curve, Convex, Pendle & 5--20\% \\
BTC Native & Babylon, Lombard, SolvBTC & 3--12\% \\
Solana & Marinade, Jito, Kamino & 5--10\% \\
TON & Tonstakers, Bemo, STON.fi & 5--15\% \\
Perps & LPX Perps LP (fee-based) & 10--30\% \\
\bottomrule
\end{tabular}
\caption{Yield strategy categories (29 strategies deployed)}
\label{tab:strategies}
\end{table}

\subsection{Yield Composition}

Bridge token holders earn layered yield:

\begin{equation}
  \text{APY}_{\text{total}} = \text{APY}_{\text{base}} + \text{APY}_{\text{restaking}} + \text{APY}_{\text{bridge}} + \text{APY}_{\text{DeFi}}
\end{equation}

For yLBTC: Babylon staking ($\sim$5\%) $+$ restaking bonus ($\sim$2\%) $+$ bridge fee share via xLUX ($\sim$1\%) $+$ DeFi yield on Lux ($\sim$5--15\%) $=$ 13--23\% total APY on BTC.

% ============================================================
\section{Shariah Compliance}
% ============================================================

\subsection{Motivation}

Islamic finance law (Shariah) prohibits \emph{riba} (interest/usury), \emph{gharar} (excessive uncertainty), and \emph{maysir} (gambling). Approximately 1.8~billion Muslims are excluded from most DeFi protocols because yields typically derive from lending interest.

\subsection{Protocol-Level Classification}

Teleport introduces the \texttt{ShariaFilter} contract:

\begin{definition}[Compliance Status]
Each yield strategy is classified as:
\begin{itemize}[nosep]
  \item \textbf{Halal}: Fee-based yield (DEX fees, bridge fees, validator staking, LP provision)
  \item \textbf{Haram}: Interest-based yield (Aave, Compound, MakerDAO DSR)
  \item \textbf{Conditional}: Requires Shariah Advisory Board (SAB) review
  \item \textbf{Under Review}: Pending classification
\end{itemize}
\end{definition}

When \texttt{shariahMode} is enabled on an OmnichainRouter, only Halal-classified strategies receive deposits. A SAB multisig can update classifications transparently on-chain.

\begin{table}[h]
\centering
\small
\begin{tabular}{ll}
\toprule
\textbf{Halal (Permitted)} & \textbf{Haram (Prohibited)} \\
\midrule
DEX trading fees & Lending interest \\
Bridge fees & MakerDAO DSR \\
Validator staking & Interest-bearing stables \\
LP provision & Leveraged farming \\
Perps trading fees & \\
Babylon BTC staking & \\
\bottomrule
\end{tabular}
\caption{Shariah compliance classification}
\label{tab:shariah}
\end{table}

% ============================================================
\section{Formal Security Properties}
% ============================================================

\begin{property}[Conservation]
For every bridge token $T$:
\begin{equation}
  \texttt{totalMinted}[T] \leq \texttt{totalBacking}[T]
\end{equation}
Enforced by automatic pause when backing ratio falls below 98.5\%.
\end{property}

\begin{property}[Uniqueness]
For all source chains $C_s$ and nonces $\nu$:
\begin{equation}
  |\{tx : tx \text{ processes } (C_s, \nu)\}| \leq 1
\end{equation}
Enforced by atomic nonce bitmap update within the mint/release transaction.
\end{property}

\begin{property}[Liveness]
If at least $t$ of $n$ MPC nodes are online and connected, the bridge processes new deposits within one source-chain confirmation period.
\end{property}

\begin{property}[Safety]
If fewer than $t$ MPC nodes are compromised, no unauthorized minting can occur. Formally:
\begin{equation}
  |\mathcal{A} \cap \mathcal{S}| < t \implies \Pr[\text{forge}(\sigma)] \leq \text{negl}(\lambda)
\end{equation}
where $\lambda$ is the security parameter.
\end{property}

\begin{property}[Governance Isolation]
The governor $G$ of chain $C$ cannot:
\begin{enumerate}[nosep]
  \item Call \texttt{bridgeMint} without valid MPC signature $\sigma$
  \item Modify processed nonce bitmaps
  \item Set \texttt{bridgeFeeBps} $> 100$ (1\% hard cap)
  \item Bypass the 48-hour timelock on stakeholder vault changes
\end{enumerate}
\end{property}

\begin{property}[Exit Guarantee]
Users can burn bridge tokens and receive underlying assets within 7 days under all conditions, including governor compromise, single MPC node failure, and pending signer rotation.
\end{property}

% ============================================================
\section{Supported Execution Environments}
% ============================================================

Teleport natively supports 20 execution environments through chain-specific bridge programs:

\begin{table}[h]
\centering
\scriptsize
\begin{tabular}{rlll}
\toprule
\textbf{\#} & \textbf{Runtime} & \textbf{Language} & \textbf{Sig Scheme} \\
\midrule
1 & EVM & Solidity & ECDSA (secp256k1) \\
2 & SVM & Rust (Anchor) & Ed25519 \\
3 & TVM (TON) & FunC & Ed25519 \\
4 & Sui MoveVM & Move & Ed25519 \\
5 & Aptos MoveVM & Move & Ed25519 \\
6 & CosmWasm & Rust & Ed25519 \\
7 & OP\_NET & AssemblyScript & Schnorr (Taproot) \\
8 & Bitcoin Script & FROST Taproot & Schnorr (BIP-340) \\
9 & XRPL Hooks & C (WASM) & ECDSA \\
10 & Substrate (ink!) & Rust & Ed25519/Sr25519 \\
11 & NEAR WASM & Rust & Ed25519 \\
12 & Soroban & Rust & Ed25519 \\
13 & Clarity & Clarity & secp256k1 \\
14 & Plutus & Aiken & Ed25519 \\
15 & Cairo & Cairo & ECDSA (Stark) \\
16 & ICP Canisters & Rust & t-ECDSA \\
17 & AVM & PyTeal & Ed25519 \\
18 & FuelVM & Sway & ECDSA \\
19 & Michelson & CameLIGO & Ed25519 \\
20 & TVM (TRON) & Solidity & ECDSA \\
\bottomrule
\end{tabular}
\caption{Native bridge programs by execution environment}
\label{tab:runtimes}
\end{table}

Total: 270 chain IDs registered, covering every top-200 blockchain by TVL.

% ============================================================
\section{Fee Economics}
% ============================================================

\subsection{Fee Structure}

Bridge operations incur a configurable fee $f_b$ (default 10~bps, max 100~bps):

\begin{equation}
  \text{mintAmount} = \text{depositAmount} \times \left(1 - \frac{f_b}{10{,}000}\right)
\end{equation}

Yield operations incur a performance fee $f_y$ (default 1000~bps = 10\% of yield):

\begin{equation}
  \text{yieldToHolders} = \text{totalYield} \times \left(1 - \frac{f_y}{10{,}000}\right)
\end{equation}

\subsection{Fee Distribution}

Fees are split between the sovereign stakeholder vault and protocol treasury:

\begin{equation}
  \text{toStakeholders} = \text{fee} \times \frac{f_s}{10{,}000}, \quad
  \text{toTreasury} = \text{fee} - \text{toStakeholders}
\end{equation}

where $f_s$ is the stakeholder share (default 9000~bps = 90\%).

\subsection{xLUX Flywheel}

On the Lux network, the stakeholder vault is LiquidLUX (xLUX). All protocol fees from bridge, DEX, lending, perps, and NFT AMM flow to xLUX holders:

\begin{equation}
  \text{APY}_{\text{xLUX}} = \frac{\sum_{\text{protocol}} \text{fees}_{\text{protocol}}}{\text{totalStaked}_{\text{LUX}}}
\end{equation}

This creates a positive feedback loop: more bridged assets $\to$ more yield $\to$ more fees $\to$ higher xLUX APY $\to$ more LUX staked $\to$ deeper liquidity.

% ============================================================
\section{Implementation}
% ============================================================

All source code is open-source under the MIT license:

\begin{table}[h]
\centering
\small
\begin{tabular}{ll}
\toprule
\textbf{Component} & \textbf{Repository} \\
\midrule
Bridge contracts & \texttt{luxfi/standard} \\
MPC signing & \texttt{luxfi/mpc} \\
Bridge infra & \texttt{luxfi/bridge} \\
Threshold crypto & \texttt{luxfi/threshold} \\
Cross-chain msg & \texttt{luxfi/warp} \\
Chain registry & \texttt{ChainIds.sol} \\
Yield strategies & \texttt{strategies/*.sol} \\
\bottomrule
\end{tabular}
\caption{Open-source repositories}
\label{tab:repos}
\end{table}

\subsection{Gas Costs}

Signature verification costs on EVM:

\begin{table}[h]
\centering
\small
\begin{tabular}{lrrr}
\toprule
\textbf{Operation} & \textbf{Gas} & \textbf{Cost (30 gwei)} \\
\midrule
\texttt{mintDeposit} & $\sim$85,000 & $\sim$\$4.50 \\
\texttt{burnForWithdrawal} & $\sim$65,000 & $\sim$\$3.40 \\
\texttt{updateBacking} & $\sim$45,000 & $\sim$\$2.35 \\
FROST verify (precompile) & $\sim$65,000 & $\sim$\$3.40 \\
CGGMP21 verify & $\sim$105,000 & $\sim$\$5.50 \\
\bottomrule
\end{tabular}
\caption{Gas costs for core operations (Ethereum mainnet)}
\label{tab:gas}
\end{table}

% ============================================================
\section{Related Work}
% ============================================================

\textbf{Wormhole}~\cite{wormhole} uses a guardian network with $13$-of-$19$ multisig. Teleport's $2$-of-$3$ threshold is more capital-efficient but assumes fewer, more trusted nodes.

\textbf{LayerZero}~\cite{layerzero} provides a generic messaging layer with configurable security (Oracle + Relayer). Teleport is bridge-specific, optimizing for asset transfer rather than general messaging.

\textbf{IBC}~\cite{ibc} achieves trust-minimized transfers via light client verification. Teleport's CosmWasm bridge can optionally use IBC for Cosmos-to-Cosmos transfers, falling back to MPC for non-IBC chains.

\textbf{tBTC}~\cite{tbtc} uses threshold ECDSA for Bitcoin bridging. Teleport generalizes this to all chains and adds yield-bearing tokens, sovereign governance, and Shariah compliance.

\textbf{Axelar}~\cite{axelar} provides a cross-chain communication network with delegated Proof-of-Stake. Teleport's MPC approach avoids the need for a separate consensus chain.

% ============================================================
\section{Conclusion}
% ============================================================

Lux Teleport provides the infrastructure for a truly omnichain liquidity layer: non-custodial, non-upgradeable, sovereignly governed, optionally Shariah-compliant, and natively deployed across every major blockchain execution environment.

By separating the attestation layer (MPC threshold signatures) from the governance layer (per-chain DAO), Teleport enables any blockchain to join the liquidity network while retaining full sovereignty over its economic parameters. The yield-bearing bridge token system makes bridged capital productive, and the Shariah compliance mode opens DeFi to 1.8~billion Muslims worldwide.

The protocol is fully open-source, with 270+ registered chains, 29 yield strategies, 20 native execution environment adapters, and a complete LaTeX-formalized security model.

% ============================================================
% References
% ============================================================
\begin{thebibliography}{99}
\bibitem{frost} C.~Komlo and I.~Goldberg, ``FROST: Flexible Round-Optimized Schnorr Threshold Signatures,'' in \emph{Selected Areas in Cryptography (SAC)}, 2020.

\bibitem{cggmp21} R.~Canetti, R.~Gennaro, S.~Goldfeder, N.~Makriyannis, and U.~Peled, ``UC Non-Interactive, Proactive, Threshold ECDSA with Identifiable Aborts,'' in \emph{ACM CCS}, 2020.

\bibitem{bip340} P.~Wuille, J.~Nick, and T.~Ruffing, ``BIP-340: Schnorr Signatures for secp256k1,'' Bitcoin Improvement Proposal, 2020.

\bibitem{bip341} P.~Wuille, J.~Nick, and A.J.~Towns, ``BIP-341: Taproot: SegWit version 1 spending rules,'' Bitcoin Improvement Proposal, 2020.

\bibitem{eip712} R.~Floersch and J.~Baylina, ``EIP-712: Typed Structured Data Hashing and Signing,'' Ethereum Improvement Proposal, 2017.

\bibitem{wormhole} Wormhole Foundation, ``Wormhole: A Generic Message Passing Protocol,'' Whitepaper, 2022.

\bibitem{layerzero} R.~Zarick, B.~Pellegrino, and C.~Banister, ``LayerZero: An Omnichain Interoperability Protocol,'' Whitepaper, 2022.

\bibitem{ibc} C.~Goes \emph{et al.}, ``The Interblockchain Communication Protocol: An Overview,'' Cosmos Documentation, 2020.

\bibitem{tbtc} M.~Luongo and C.~Ream, ``tBTC: A Decentralized Redeemable BTC-backed ERC-20 Token,'' Keep Network, 2020.

\bibitem{axelar} S.~Layr \emph{et al.}, ``Axelar: Connecting Web3,'' Whitepaper, 2022.

\bibitem{rekt} Rekt News, ``Leaderboard,'' \url{https://rekt.news/leaderboard/}, 2024.

\bibitem{erc4626} J.~Bebis \emph{et al.}, ``EIP-4626: Tokenized Vault Standard,'' Ethereum Improvement Proposal, 2022.
\end{thebibliography}

\end{document}
